% Ad Familiares 6.21 --- headnote
% ad-familiares-06-21, around April 45 BC, Rome

\textit{Cicero to C.~Toranius, written at Rome around April of
45 BC (Perseus: \textit{Romae circ.~m.~Apr.~a.~709 (45)}). Toranius
was a former aedile-colleague of Octavius's father and another
Pompeian who had not been restored after Pharsalus; he was still
living away from Italy. The Spanish war was at its hinge --- the
letter is written, on Cicero's own account, with the impression
that either the end is at hand or has already come (the battle of
Munda was fought on 17 March, but the news had not yet reached
Rome) --- and Cicero takes the moment to reach across to a man
who had stood at his shoulder in the council-of-war debates of
49 BC against Domitius Ahenobarbus and the Lentuli.}

\textit{The argument is the consolatory commonplace of the
period in its philosophical form: those who saw clearly what war
would bring, and were called timid for saying so, have at any
rate this comfort --- that they judged the case rightly and
truly. Whatever follows can be borne with moderation, since
death is the end of all things and the conscience of having
acted for the commonwealth's dignity remains. The letter then
turns from himself to Toranius and closes with the assurance
of continued service to him, to his safety, and to his children.
The register sits closer to the Torquatus consolations
(\textit{Fam.}~6.1--4) than to the warmer family-letter manner
of the Lepta correspondence, but Cicero is writing to a more
distant friend and the tone is correspondingly measured.}
